% ======================= 第 十 章 =======================
\chapter{名字的层累}
\chaptersubtitle{ГЛАВА ДЕСЯТАЯ}{СЛОИСТОСТЬ ИМЁН}
\epigraphbox{“名者，实之宾也。”}{《庄子·逍遥游》}

\kaishou{本章只做一件事：指认。传世文献里是四个人的，在档案里是同一批人；至于这些名字怎么变的，那是另一本书的题目。}

\noindent
前九章讲完了这个文明的一生。剩下的事是它怎么变成了神话。

这一章的方法很笨：把传世文献里的名号与档案里的职掌一个一个对上，能对上的写出来，对不上的注明对不上。

\section{四个名号与四项职掌}

\begin{table}[htbp]
\centering
\caption{四个名号与档案里的对应}
\label{tab:four}
\small
\begin{tabular}{@{}llp{66mm}@{}}
\toprule
名号 & 出处 & 档案里的对应\\
\midrule
嫦娥 & 《淮南子·览冥训》 & 驻地球主任，前 2050 年下地\\
西王母 & 《山海经·西山经》 & 驻地球主任的上级，即后来的八意永琳\\
后羿 & 《左传》襄公四年 & 有穷氏首领，前 2000 年代夏\\
纯狐 & 《左传》昭公二十八年 & 后羿之妻、伯封之母，有穷氏土著\\
\bottomrule
\end{tabular}
\tabsource{出处一栏只给最早见的文献，不列全部。对应一栏为本书所拟，两处都不见于同名文本。}
\end{table}

表 \ref{tab:four} 的四个名号里，前三个在传世文献里各有完整的故事，第四个只有几句。

本书要说的是：这四个人在档案里分别对应四个职掌，而这四个职掌之间的关系是一套行政关系，不是一套神话关系。

\section{姮娥到嫦娥的避讳}

嫦娥这个名字，在传世文献里最早的写法是「姮娥」。《淮南子·览冥训》记「羿请不死之药于西王母，姮娥窃以奔月」。

改作「嫦」是避讳。汉文帝名恒，恒与姮同音，改字避之。此后一律作嫦娥。

这一条没有争议，历来注家都这么说。本书要补的是最后一句：此后一律作嫦娥，连档案也跟着改了。卷一所辑的残简里，她只出现过一次，记作「嫦娥氏」，而那片简的年代在前二二一年之后，属后世题记。

也就是说，这个人的本名在材料里根本没有留下来。

\section{相隔八百年的两个形象}

《山海经·西山经》记西王母「其状如人，豹尾虎齿而善啸，蓬发戴胜」。

六朝以后的西王母是另一副样子：《汉武帝内传》里她是一位绝世女仙，年三十许，容颜绝世，修短得中。

两个形象之间隔了至少八百年。传统的解释是神话演变，从图腾神变成女仙。

本书的解释更简单：这两个形象不是同一个人。前一个是对某位外来者的记忆，后一个是对那段记忆的改写。

改写的方法，见《记忆消除场》那本书，此处不重复。本书只指出一件事：改写发生在中原，不在月面。留在月面的档案里，这个人的职掌写得清清楚楚，是一个管地面对接的官员，与豹尾虎齿没有关系。

\section{《左传》中的两处记载}

后羿在《左传》里是有穷氏的首领，在《淮南子》里是射日的英雄。两个身份之间隔着一次改历，见第五章。

纯狐在《左传》昭公二十八年只出现一次：后羿「灭伯封」，伯封之母曰玄妻，有仍氏之女。后世注家把玄妻与纯狐氏并为一事，这个合并是注家的，不是原文的。

本书取合并说，理由只有一条：若不合并，纯狐这个名字在材料里就没有来源。合并之后，她的身份对得上一处：一位有穷氏的土著女性，手上有旧历抄本，幼年见过月人遗迹的刻痕。

一个不识字的人，凭这两样东西，用了那台机器。第六章已经说过，这里不重复。

\section{神话化的条件}

四个人的名字都进了神话，而这座城本身没有。

这是本书最不愿意碰、又不得不碰的一处。原因不难找：神话是给活人用的，城是给历史用的。一座在月亮上的城，对写书的人没有用处；一个能赐药的女仙，对求长生的人有用处。

有用处的东西被反复讲，反复讲的东西变形。没有用处的东西只剩下一个地名，被顺手写进志怪里，随后连地名也散了。

这一条与本书的总题目是同构的：他们把「变化」删掉了，于是两万年不动；不动的东西没有人讲，于是它变成神话，也变成两副面孔。

\section{本章结论}

本章说了四件事。

嫦娥与姮娥是同一个人，改名是避讳，本名没有留下来。

西王母有两个形象，前后相隔八百年，本书认为后一个是改写。

后羿与纯狐在《左传》里只有几句，本书把它们与改历和那台机器对应起来。

这座城没有进神话，它的人进了神话。

本章的自我拆台：这一章的对上，全部是职掌与情节的对应，不是名字的对应。档案里没有「嫦娥」二字，也没有「西王母」三字，只有职掌。从职掌推名字，中间隔着一层，这一层本书填不上。

笔者把这层留着，没有用「显然」「无疑」之类的词糊过去。理由是：这四个名字在通俗读物里另有来历，本书与那是两回事，不必互相印证。
